\documentclass{article}
\usepackage[utf8]{inputenc}
\usepackage[T1]{fontenc}
\usepackage{geometry}
\usepackage{graphicx}
\usepackage{float}
\usepackage{hyperref}
\usepackage[english,indonesian]{babel}

\title{Optimasi Pemrosesan Template LaTeX Menggunakan Pendekatan Berbasis AI}
\author{Academic Support Agent}
\date{\today}

\begin{document}
\maketitle

\begin{abstract}

\end{abstract}

\chapter{PENDAHULUAN}


\chapter{TINJAUAN PUSTAKA}


% Original content (commented out):
% \documentclass[12pt]{article}
\usepackage[margin=1in]{geometry}% Change the margins here if you wish.
\setlength{\parskip}{5pt} % This sets the distance between paragraphs
\usepackage[utf8]{inputenc}
\usepackage{amsmath}
\usepackage{mathtools}
\usepackage[framed,numbered,autolinebreaks]{mcode}
\usepackage{physics}
\usepackage{graphicx}
\usepackage{siunitx}


\title{ME C134: Prelab 2}
\author{Jordan Suchard, Wednesday 2--5}
\date{\today}

\begin{document}

\maketitle

\section{Purpose}

\section{Prelab}

\section{Lab}

\appendix

\section{Appendix}



\end{document}


\end{document}